\chapter{待解决的研究问题}

当前真正可能改变顶层 cognition ownership 的问题已经明显减少。CognitiveSchema 归属已收敛为 Memory domain 内的一等对象族；Social Cognition 的方向性关系图和 LanguageConvention 也已形成目标语义。以下问题继续保持 OPEN，避免用实现便利提前冻结。

\section{宏观语义与领域所有权}

\begin{enumerate}
\item \textbf{Self / Cognitive Seed constitution：} Seed 中哪些内容只是可演化的初始条件，哪些属于不能由普通 cognition 自我覆盖的 constitutional constraint？Self Model、Persona/dispositions、values/preferences 与 narrative identity 如何在同一 Self owner 中保持不同修订和冲突语义？管理员修改 Seed 时，是修改初始化来源、产生新的 constitution revision，还是提交 Self proposal？
\item \textbf{Motivation 范围与生命周期：} values、needs、interests 分别属于 Self 还是 Motivation？Goal 的激活、优先级、冲突、满足、暂停、放弃和过期需要哪些最小 durable semantics？Runtime priority 与长期 Motivation 如何分离？
\end{enumerate}

Subject fork/clone 的原则已经确定为新 SubjectId + inherited immutable lineage + branch-local mutable heads；具体 copy-on-write、purge lineage 和导出/迁移合同属于实现/产品层，不再作为顶层 ontology blocker。

\section{Social Cognition 的规格问题}

Social 顶层 owner 和两个首批对象族已经确定，剩余问题应在 Social Spec 中收敛，而不是重新打开领域边界：

\begin{itemize}
\item relation kind 的最小 canonical set，以及哪些关系允许 ordinal degree；
\item degree 的 operational semantics，避免跨 relation kind 的无定义算术比较；
\item Person/Dyad/Group/Community/Channel 等 LanguageConvention scope 的最小编码；
\item 显式说明、重复成功指代、repair/acknowledgement 等 evidence 如何进入 formation proposal；
\item competing scoped meanings、relation contradictions、split/merge 和 identity rebind 的 revision contract；
\item social graph 的 Serving projection 如何处理 fan-out、hub、方向性和 query-specific inverse view。
\end{itemize}

这些问题不会阻塞 Memory Reference Profile。

\section{记忆、运行时与检索实验}

以下问题已经有稳定 semantic seam，但默认算法/参数需要真实工作负载验证：

\begin{itemize}
\item AUTO accessibility：比较无硬衰减、time-only、time+typed-use 以及更复杂候选模型；
\item multi-view aggregation：best-view、capped multi-view 及其 recall/cost；
\item topology utility：在相同 hard filters、图、种子和预算下比较 no-topology、typed PPR、state-faithful Wave 等；
\item EPA：修正 weighted covariance 后再评价是否带来弱线索检索收益；
\item Residual：与普通 dense、多查询和模型分解比较真实 retrieval utility；
\item reranking：确定在 PC-local workload 中何时收益足以抵消模型成本；
\item dense backend：按 filtered recall、multi-view、RAM/disk、build/restart、backup/restore 和部署复杂度比较 pgvector/Qdrant 等候选。
\end{itemize}

算法实验不能改变 Authority identity、Evidence、Revision、UseEvent、Suppression/Purge 等基础语义。

\section{Episode、Journal 与 Offline Cognition 的实现合同}

Episode/Journal 的对象性质已明确，仍需在实现前写清：

\begin{itemize}
\item Episode boundary proposal 的输入信号、split/merge/re-segment 与历史引用解析；
\item Journal source revision、复核状态、重写边界和 retrieval role；
\item Heptalogos maintenance opportunity/lease 的资源、取消和 deadline 语义；
\item Dream seed selection、novelty/quality stop、循环抑制、counterevidence search 和 proposal routing；
\item Hypothesis 的 stable identity、输入引用、producer signature、uncertainty 和 owner decision receipt。
\end{itemize}

这些合同属于后续 Memory/offline cognition wave，不作为 10 月 27 日 Memory Reference Profile 的必要条件。

\section{可靠性与跨域合同}

需要落到 Spec/transaction/test 的问题包括：

\begin{itemize}
\item UseEvent request/event idempotency 与 timeout-after-commit；
\item projection outbox/watermark、stale result revalidation 和 rebuild failure；
\item purge 对 Authority、Serving、cache/context、material、backup/export retention 的一致语义；
\item 来源访问权撤销和 derived material 的 invalidation；
\item Heptalogos 与 Nous Wave 的 Subject provisioning、Observation、Context、Goal/Commitment、Use/Outcome、Maintenance handoff；
\item 不使用跨库分布式 transaction 时的可靠事件交接、receipt 和重试。
\end{itemize}

\section{从研究到执行规格}

当前最大工程缺口已经不是继续增加宏观概念，而是把已收敛的设计下降为可执行合同：

\begin{center}
owner/object semantics $\rightarrow$ operations $\rightarrow$ Proto/SQL/types $\rightarrow$ transaction/CAS/idempotency $\rightarrow$ projection invalidation $\rightarrow$ acceptance tests $\rightarrow$ implementation Plan
\end{center}

Memory 是第一套走完整条链路的 Reference Profile；它完成后，Social、Motivation 和 Self 应复用同一 Evidence/Revision/Proposal/Use 基础，而不是各自重新发明 cognition infrastructure。